% =============================================================
%  第一单元  100 以内数加与减（二）
% =============================================================
\Chapter{第一关\quad 一串小棒走天下}
\noindent{\small\kaishu\color{tipc}（课本第一单元：100 以内加与减（二），P2--P17）}

\begin{blunder}
开学第一周，超市结账 48 块，爸爸给了 100，收银员找给 52。
爸爸心里“咦”了一下——怎么才找这么点？回家一算，是爸爸自己心算错了。
更尴尬的是妈妈在旁边补了一句：“一百减四十八，不就是五十二吗。”
\end{blunder}

\begin{gamble}{课前协商：五道“100 减两位数”}
\begin{itemize}
  \item 两人轮流出题：爸爸出五道给小豆角算，小豆角也出五道考爸爸。
  \item \textbf{小豆角当裁判}：爸爸做题时，小豆角负责挑错，抓到一处记一颗星。\waterdrop
  \item 三局两胜，打完当场兑现。
\end{itemize}
\negolines
\end{gamble}

\begin{bonus}
\begin{itemize}
  \item \textbf{电梯挑战}：进电梯让小豆角闭上眼，猜“从 1 楼到 12 楼经过几层”。
        到了睁眼看楼层数字。猜对一次记一颗星。
  \item \textbf{存钱罐对账}：每周数一次零花钱，算“这周比上周多了多少”。
        多出来的钱小豆角可以自己决定怎么花——这比十道应用题管用。
  \item \textbf{找零游戏}：超市结账前让小豆角先估个总数，结完对一对。
        估差在 5 块以内，这周小豆角挑一次零食。
\end{itemize}
\end{bonus}

\begin{unitinfo}
\textbf{课本 P2--P17}（外加“数学好玩·猜数游戏”P18--P19）\quad
\textbf{6 次课，每次 20 分钟} \\[4pt]
\textbf{这一关到底在干嘛}：就是把“10 根捆成 1 捆”和“拆开 1 捆当 10 根”
这两个动作，反复做到手上。\\
\textbf{为什么它最要紧}：后面乘、除、口诀，全都踩在这两个动作上。
这两个动作要是只“背下来”，到三年级学多位数、四年级学小数，会整片塌掉。
\textbf{宁可慢，也必须亲手捆、亲手拆。}
\end{unitinfo}

\subsection{开局先聊个事：这些题生活里长什么样}

爸爸今天带小豆角去超市，东西 48 块，爸爸给 100，收银员找爸爸多少？——这就是 $100-48$。

电梯在 1 楼，小豆角去 12 楼的同学家玩，中间经过几层？
——这就是“差几层”。

小豆角跳绳跳了 48 个，隔壁小朋友比小豆角多跳 25 个，人家跳多少？
——这就是“比谁多”。

课本把这些都编成了小故事（摘苹果、收玉米、跳绳），
本质是同一件事：\textbf{小豆角要是能看着东西数出来，就能算出来。}

\subsection{任务 1：把“满十进一”捆出来}

\begin{lesson}{1}{20}
\textbf{今天要拿下的}：让小豆角自己把 10 根散的小棒捆成一捆，
并且说出“这一捆是从散的那堆里来的”。

\begin{materials}
小棒 60 根、皮筋 10 根、一张白纸、一支笔。
\end{materials}

\begin{enumerate}
  \item \textbf{摆}：桌上摆 35 根——3 捆（每捆 10 根）＋ 5 根散的。
        让小豆角数一遍，说清楚“3 个十和 5 个一”。
  \item \textbf{加}：旁边再放 8 根散的。问：“现在散的一共有几根？”
        （$5+8=13$ 根）
  \item \textbf{等小豆角捆}：把皮筋递过去，只说一句：“散的有点多了，小豆角想想办法。”
        \textbf{全程不要给提示，不说话就好了}\\
        小豆角要是想不到，就提示一句：“10 根散的一捆，是什么？”
        \textbf{不要直接说“捆起来”。}
  \item \textbf{数}：捆好以后数——4 捆 ＋ 3 根散 $=43$。
  \item \textbf{写}：让小豆角自己在纸上写 $35+8=43$。
  \item \textbf{问那句关键的话}：“本来是 3 捆，怎么变成 4 捆了？”\\
        小豆角答“散的凑够 10 根，捆了一捆”——\textbf{好，这就是“满十进一”}。
        让小豆角自己说出这句话，别替小豆角说。
  \item \textbf{换题再走一遍}：$26+19$。19 得先自己捆成 1 捆＋9 根，
        然后 2 捆＋1 捆，6 根＋9 根 $=15$ 根，再捆一捆 →
        4 捆＋5 根 $=45$。
\end{enumerate}

\begin{saybox}
“这一捆是从哪里来的？”\\
“3 捆怎么变成 4 捆的？”\\
“散的满 10 根了，小豆角打算怎么办？”\\
“小豆角刚才做的这个动作，不用小棒的话，怎么记下来？”
\end{saybox}

\begin{pitbox}
\textbf{最常见的卡点}：小豆角算出“13 根散的”，然后就把 13 根散着摆在那儿，
直接数出“3 捆 + 13 根”$=43$——\textbf{答案对，但“捆”这个动作没做}。\\
\textbf{怎么办}：别纠答案（答案没错），追一句：“13 根散的，能不能变成 1 捆几根？”
让小豆角把那一捆补上。\textbf{这一课全部的价值，就在这个动作里。}
\end{pitbox}

\begin{passbox}
\begin{itemize}
  \item “摆 $\to$ 加 $\to$ 捆 $\to$ 数”四步小豆角自己走完，爸爸一根手指都不动。
  \item 小豆角能用自己的话讲清“为什么十位多了一”。
  \item 换成 $38+9$ 或 $17+25$，小豆角照样能走一遍。
\end{itemize}
\end{passbox}
\end{lesson}

\subsection{任务 2：竖式那三条规矩，和“当医生”}

\begin{lesson}{2}{20}
\textbf{今天要拿下的}：把昨天手上捆的动作，翻译成竖式上的三个规矩。

\begin{materials}
方格竖式纸（附录 B）、笔。有条件用白板，能擦了重写。
\end{materials}

\begin{enumerate}
  \item \textbf{先用手再做一遍} $38+9$，三十秒，确认动作还记得。
  \item \textbf{搬到纸上}：拿方格竖式纸，一位一格。
        \begin{center}
        \begin{tabular}{@{}c@{\hspace{1.6cm}}c@{}}
          \begin{shushicell}
            & 3 & 8\\
          + &   & 9\\ \midrule
            & 4 & 7
          \end{shushicell}
          &
          \begin{shushicell}
            & 2 & 6\\
          + & 1 & 9\\ \midrule
            & 4 & 5
          \end{shushicell}
        \end{tabular}
        \end{center}
        \textbf{边写边问}：“个位 8+9=17，格子只有一个，17 放哪？”
        （写 7，把 1 记在十位上面，用小字写。）\\
        “那个小小的 1，是 1 个一还是 1 个十？”——1 个十，所以它站十位。
  \item \textbf{连做四道}：$38+9$、$27+24$、$35+15$、$9+38$。
        每道都要求\textbf{先把进位小 1 写出来}，再算十位。
  \item \textbf{让小豆角当医生}：下面三个病人，找病因、改正。
        \begin{center}
        \begin{tabular}{@{}c@{\hspace{1.1cm}}c@{\hspace{1.1cm}}c@{}}
          \begin{shushicell}
            & 4 & 7\\
          + &   & 6\\ \midrule
            & 4 & 3
          \end{shushicell}
          &
          \begin{shushicell}
            & 3 & 6\\
          + & 2 & 9\\ \midrule
            & 5 & 5
          \end{shushicell}
          &
          \begin{shushicell}
            & 5 & 2\\
          + & 1 & 7\\ \midrule
            & 6 & 9
          \end{shushicell}
        \end{tabular}
        \end{center}
        答案：第一题 53（个位 13 只写了 3，把进的 1 丢了）；
        第二题 65（十位 3+2 忘了加进的 1）；第三题 69——\textbf{这个是对的，故意放个假病人}。
  \item \textbf{收三句话}（让小豆角自己说，爸爸写在纸上贴起来）：
        ① 个位对个位，十位对十位；② 个位满十，在十位上面写个小 1；
        ③ 十位相加时，别忘那个 1。
\end{enumerate}

\begin{saybox}
“小豆角把进的 1 放哪了？为什么放十位，不放个位？”\\
“这道题小豆角是医生，小豆角说它得了什么病？”（比“小豆角又算错了”印象更深）\\
“十位要是忘了加 1，得数会差多少？”（差 10，让小豆角自己发现）
\end{saybox}

\begin{passbox}
\begin{itemize}
  \item 5 道进位加全对，而且\textbf{每道都有进位小 1 的痕迹}。
  \item 前两个“病人”的病因小豆角自己找出来了。
\end{itemize}
\end{passbox}
\end{lesson}

\subsection{任务 3：把“退一当十”拆出来}

\begin{lesson}{3}{20}
\textbf{今天要拿下的}：小豆角会亲手拆开一捆，并说出“十位要少 1”。

\begin{materials}
小棒、皮筋、笔。
\end{materials}

\begin{enumerate}
  \item \textbf{摆} 52：5 捆＋2 根。要拿走 7 根。
  \item \textbf{先卡住小豆角}：让小豆角拿 7 根。散的只有 2 根，\textbf{拿不出来}。
        停在这儿，问：“怎么办？”
  \item \textbf{拆}：把 1 捆拆开，变 10 根散的 → 现在散的 12 根，捆的剩 4 捆。
        拿走 7 根，散剩 5 根。结果 4 捆＋5 根 $=45$。
  \item \textbf{搬到纸上}：
        \begin{center}
        \begin{shushicell}
          & 5 & 2\\
        - & 3 & 7\\ \midrule
          & 1 & 5
        \end{shushicell}
        \end{center}
        “个位 2 减 7 不够，从十位借 1 当 10，$12-7=5$；
        十位 5 已经借走 1，只剩 4，$4-3=1$。”
  \item \textbf{教一个招：退位点}。借过位的十位，在上面点一个小点。
        \textbf{看到点就先减 1}——这是防“借了忘还”的唯一办法。
        让小豆角以后每道退位题都点。
  \item \textbf{连做四道}：$52-37$、$53-26$、$32-3$、$43-14$。
  \item \textbf{当医生}：
        \begin{center}
        \begin{tabular}{@{}c@{\hspace{1.6cm}}c@{}}
          \begin{shushicell}
            & 5 & 2\\
          - &   & 8\\ \midrule
            & 5 & 4
          \end{shushicell}
          &
          \begin{shushicell}
            & 4 & 3\\
          - & 1 & 4\\ \midrule
            & 3 & 1
          \end{shushicell}
        \end{tabular}
        \end{center}
        答案：第一题 44（把 $8-2$ 倒着减了）；第二题 29（十位忘了已经借走 1）。
\end{enumerate}

\begin{saybox}
“小豆角借来的那一捆，是几个十？那十位是不是要少 1？”\\
“个位不够减，能不能拿大的减小的？”（不能——这条小豆角最想干）\\
“小豆角十位上那个点，什么时候点的？为什么点它？”
\end{saybox}

\begin{pitbox}
\textbf{三个高频错误，盯着看：}
\begin{enumerate}[label=\arabic*、]
  \item \textbf{倒着减}：$52-8$ 算成 54。\textbf{一句话}：“不够减就先借，不许倒着减。”
  \item \textbf{借了不还}：$43-14$ 算成 39。\textbf{一句话}：“点上点的那一位，先减 1 再算。”
  \item \textbf{借来的当成 1}：以为借来的是 1 而不是 10。
        \textbf{一句话}：回小棒，“小豆角拆开那一捆，是 1 根还是 10 根？”
\end{enumerate}
\end{pitbox}

\begin{passbox}
\begin{itemize}
  \item 用小棒独立做出 $52-37$，包括“拆捆”那一下。
  \item 5 道退位竖式全对，每题都有退位点。
  \item 小豆角能讲清“倒着减”为什么不对。
\end{itemize}
\end{passbox}
\end{lesson}

\subsection{任务 4：100 减两位数（难度较大·重点关注）}

\begin{lesson}{4}{25}
\textbf{今天要拿下的}：十位是 0 借不到的时候，知道“连着借两次”。

\begin{materials}
计数器（这一课\textbf{必须用}，比小棒直观得多）。
\end{materials}

\begin{enumerate}
  \item \textbf{拨} 100：只在百位拨 1 颗珠。
  \item \textbf{卡住小豆角}：要减 48。个位 0 颗，减不了 8；十位也 0 颗，借不到。
        停住，问：“怎么办？”
  \item \textbf{借两次}（让小豆角自己动手挪珠子）：
        \begin{itemize}
          \item 第一步：从百位借 1 颗放十位 → 十位变成 10 颗（$=10$ 个十）。
          \item 第二步：从十位借 1 颗放个位 → 个位 10 颗，十位剩 9 颗。
        \end{itemize}
        现在：百位 0、十位 9、个位 10。
  \item \textbf{减}：个位 $10-8=2$；十位 $9-4=5$ → 52。
  \item \textbf{写竖式}：
        \begin{center}
        \begin{shushicell}
          & 1 & 0 & 0\\
        - &   & 4 & 8\\ \midrule
          &   & 5 & 2
        \end{shushicell}
        \end{center}
        让小豆角自己：百位 1 上打点、十位 0 改成 9、个位 0 改成 10。
        \textbf{三个记号都要画出来}，画全了才算过。
  \item \textbf{连做}：$100-39$、$100-24$、$70-27$、$80-55$。
        （$70-27$ 只借一次：十位 7 借 1 给个位，个位变 10，十位剩 6，$6-2=4$ → 43。）
  \item \textbf{找一个技巧}：不会算的时候，\textbf{先估一个数}。
        $100-48$ 大概是多少？50 上下。算出来 52——就在 50 附近，\textbf{这说明小豆角没算错方向}。
        养成“先估后算”的习惯，比多做题管用。
\end{enumerate}

\begin{saybox}
“十位是 0，借不到，小豆角想向谁借？”\\
“从百位借来的 1，到了十位变成几了？”（10 个十）\\
“十位有 10 了，个位还是 0，怎么办？”（再从十位借 1）\\
“小豆角估一下，$100-48$ 大概多少？”（先估再算）
\end{saybox}

\begin{pitbox}
\textbf{两种典型错法}：
\begin{enumerate}[label=\arabic*、]
  \item 算出 62 或 68——十位当成 10 了，借了没还。\\
\textbf{解法}：让小豆角数计数器上十位到底剩几颗（9 颗），眼见为实。
  \item 直接放弃，说“不会”。\\
\textbf{解法}：一定先上计数器。纸上只是记录，动手在前面。
\end{enumerate}
\textbf{这一课别赶。}25 分钟只弄明白 $100-48$ 一道，也值。
\end{pitbox}

\begin{passbox}
\begin{itemize}
  \item 计数器上独立做出 $100-48$，并说出“借了两次”。
  \item 竖式上三个记号齐全（百位点、十位 0 改 9、个位 0 改 10）。
  \item $100-39$、$70-27$ 独立做对。
\end{itemize}
\end{passbox}
\end{lesson}

\subsection{任务 5：先画图，再列式}

\begin{lesson}{5}{20}
\textbf{今天要拿下的}：遇到“比谁多／少”，\textbf{第一反应是画一条线}，不是猜加减。

\begin{materials}
纸、笔、尺子。
\end{materials}

\begin{enumerate}
  \item \textbf{课本 P10 跳绳}：小红跳了 48 次；小亮比小红少跳 9 次；
        小东比小红多跳 25 次。问小亮、小东各跳多少。
  \item \textbf{让小豆角圈信息}：三个数字，加“多”“少”两个字，全圈出来。
  \item \textbf{画线段}（慢一点，这一步是重点）：
        \begin{center}
        \begin{tikzpicture}[x=0.14cm,y=0.95cm,font=\small]
          \draw[mainc,thick] (0,2) -- (48,2);
          \node[left] at (0,2) {小红};
          \node[above] at (24,2) {48 次};
          \draw[pitc,thick] (0,1) -- (39,1);
          \node[left] at (0,1) {小亮};
          \node[above] at (19.5,1) {比小红少 9，所以短一点};
          \draw[sayc,thick] (0,0) -- (73,0);
          \node[left] at (0,0) {小东};
          \node[above] at (60,0) {比小红多 25，所以长一点};
          \draw[dashed,gray] (48,2) -- (48,0.15);
        \end{tikzpicture}
        \end{center}
  \item \textbf{从图上读数}：小亮 $=48-9=39$（次）；小东 $=48+25=73$（次）。
  \item \textbf{反过来问一次}：小红想跳到 100 次，还要跳多少？$100-48=52$（次）。
  \item \textbf{换个场景再来一遍}（课本 P12 图书馆三层）：
        第一层 23 本，第二层比第一层多 18 本，第三层比第一层少 9 本。
        画三条线 $\to 23+18=41$、$23-9=14$。
  \item \textbf{让小豆角总结一句话}：“‘比谁多’就在谁身上加一段；‘比谁少’就减一段。
        \textbf{所以第一步是找出‘谁’是标准。}”
\end{enumerate}

\begin{saybox}
“‘比小红少’，小豆角是拿谁当标准？”\\
“小豆角能用手指在图上指出‘多出来的 25 次’是哪一段吗？”\\
“如果不画图，小豆角能保证不把加法写成减法吗？”（让小豆角自己承认画图有用）
\end{saybox}

\begin{pitbox}
\textbf{最容易翻车的地方}：小豆角看到“少 9”就减、看到“多 25”就加，形式上对了。
可换成“小亮跳了 39 次，比小红少 9 次，小红跳多少？”小豆角立刻翻车——
因为\textbf{标准换了人}。\\
\textbf{解法}：每道题先问一句“\textbf{谁是标准}”，把标准那个人画在最上面。
位置对了，加减就不会错。
\end{pitbox}

\begin{passbox}
\begin{itemize}
  \item 三条线段小豆角自己画出来（长短、标注、对齐都对）。
  \item 小豆角能看着自己的图讲“为什么这里要加、那里要减”。
  \item 反向题（知道差和一个量，求另一个）也能做。
\end{itemize}
\end{passbox}
\end{lesson}

\subsection{任务 6：闯关 + 猜数游戏}

\begin{lesson}{6}{20}
\textbf{今天要拿下的}：前五课验一遍；顺便学一个“聪明地问”的办法。

\begin{materials}
百数表（附录 B）、纸、笔。
\end{materials}

\begin{enumerate}
  \item \textbf{先做 5 分钟快检}（见本章末），错的记下来。
  \item \textbf{猜数游戏}（课本 P18--P19）：
        \begin{playbox}{猜数（只能答“对”或“不对”）}
        ① 爸爸心里想一个 1--100 的数，小豆角来猜。小豆角问，爸爸只能说“对”或“不对”。\\
        ② 数一数小豆角问了几次猜中。\\
        ③ 再来一次，第二轮爸爸示范一句：“比 50 大吗？”然后每次都问“比中间大吗”，
        让小豆角看着范围一次砍一半：$100\to50\to25\to12\to6\to3\to1$，最多 7 次。\\
        ④ 第三轮换小豆角当“想数的人”，让小豆角体会“一半一半地排除”为什么快。
        \end{playbox}
  \item \textbf{收一句话}：“每次问‘比中间大吗’，一下就去掉一半可能，
        所以 7 次一定猜得到。”——\textbf{不用讲“二分法”这个词}，二年级只要体会。
  \item \textbf{收尾}：错的题抄进错题本（只写错因一句话），贴今天的星星。
\end{enumerate}

\begin{passbox}
\begin{itemize}
  \item 快检 10 题至少对 8 题，错的能说出错在哪一步。
  \item 猜数能在 7 次以内猜中 1--100 的任意数。
\end{itemize}
\end{passbox}
\end{lesson}

\subsection{翻车现场（对着看）}

\begin{center}
\small
\begin{tabularx}{\linewidth}{@{}p{3.6cm}p{4.0cm}X@{}}
\toprule
\textbf{小豆角写出来的样子} & \textbf{其实是想错了什么} & \textbf{爸爸回一句什么} \\
\midrule
$28+35=513$ & 不知道进位是加到\textbf{十位}上，把 1 挤在个位后面 &
“进的 1 是 1 个\textbf{十}，要住十位。” \\
\addlinespace[2pt]
$52-8=54$ & 个位不够减，倒着减了 & “不够减先借，不许倒着减。” \\
\addlinespace[2pt]
$43-14=39$ & 十位忘了已经借走 1 & “点上点的那位，先减 1。” \\
\addlinespace[2pt]
$100-48$ 空着 & 十位是 0，借不到就卡死 & 先上计数器，把“借”做出来。 \\
\addlinespace[2pt]
个位十位上下错开 & 没有格子管着，靠眼睛估 & 整章只用方格竖式纸。 \\
\addlinespace[2pt]
“比……多”做成减法 & 语言和运算没接上 & 一律先画线段，标出谁是标准。 \\
\addlinespace[2pt]
算式对，答语不写 & 没有“写完”的习惯 & 最后一行必须写“答：……”。 \\
\bottomrule
\end{tabularx}
\end{center}

\subsection{抽查：随机开枪（10 题）}

\begin{enumerate}[label=\arabic*.,itemsep=4pt]
  \item $35+8=$ \blank \qquad 2. $47+9=$ \blank \qquad 3. $26+19=$ \blank
  \item $52-7=$ \blank \qquad 5. $52-37=$ \blank \qquad 6. $100-48=$ \blank
  \item $70-27=$ \blank \qquad 8. $38+9=$ \blank \qquad 9. $60-41=$ \blank
  \item $48+25=$ \blank
\end{enumerate}

\begin{ansbox}
1. 43\quad 2. 56\quad 3. 45\quad 4. 45\quad 5. 15\quad
6. 52\quad 7. 43\quad 8. 47\quad 9. 19\quad 10. 73\\
\textbf{不到 7 题对}：回任务 2、任务 3再来一遍，别急着进第二单元。
\end{ansbox}

\subsection{平时生活里怎么练（不占上课时间）}

\begin{itemize}
  \item \textbf{超市算钱}：让小豆角算“两样一共多少钱”，钱数压在 100 以内。
        再给小豆角 100 块，让小豆角算该找多少——$100-$ 两位数的真身就在这儿。
  \item \textbf{电梯}：从 1 楼到 12 楼经过几层？再下到 3 楼，往下走几层？
  \item \textbf{扑克牌凑 100}：抽两张（A 当 1，JQK 去掉），看能不能凑成整十；
        熟了以后抽三张，先凑十再加。
  \item \textbf{存钱罐}：每周数一次零花钱，让小豆角算“这周比上周多了多少”。
  \item \textbf{每天 3 分钟口算卡}：抽 10 张，只报答案不写。连续五天全对，加一颗星。
\end{itemize}
